\section{从一个方向到直接映射}
\label{sec:method-comparison}

\subsection{共同坐标与几何差异}

三类方法可以放进同一比较坐标，但不构成理论包含链。图~\ref{fig:method-geometry} 从左到右画出单方向投影、SOM 多方向投影和模块输出的局部邻域映射。横向排列只表示决策变量发生变化：$r$、有序方向组 $(r_1,\ldots,r_k)$、以及上游矩阵 $W$（本地适配则为因子 $A,B$）。它不表示后一类方法必然包含前一类方法，也不表示拒答抑制效果随面板顺序递增。

\begin{figure*}[t]
  \centering
  \resizebox{0.97\textwidth}{!}{%
  \begin{tikzpicture}[
    font=\sffamily\scriptsize,
    >=Latex,
    panel/.style={draw=black!45, rounded corners=2pt, fill=black!1},
    good/.style={circle, draw=blue!70!black, fill=blue!18, inner sep=1.5pt},
    bad/.style={circle, draw=red!70!black, fill=red!18, inner sep=1.5pt},
    proto/.style={circle, draw=orange!80!black, fill=orange!22, inner sep=1.6pt},
    mapped/.style={rectangle, draw=green!50!black, fill=green!20, inner sep=1.5pt},
    flow/.style={-{Latex[length=2.2mm]}, line width=0.75pt},
    thinflow/.style={-{Latex[length=1.7mm]}, line width=0.55pt}
  ]

  % Panel A: single direction
  \begin{scope}[xshift=0cm]
    \draw[panel] (0,0) rectangle (6.6,4.65);
    \node[font=\sffamily\bfseries\small] at (3.3,4.34) {(a) 单方向投影};
    \node[good] at (1.15,1.20) {};
    \node[good] at (1.55,1.62) {};
    \node[good] at (1.85,0.93) {};
    \node[good] at (2.15,1.42) {};
    \node[bad] at (4.25,2.88) {};
    \node[bad] at (4.65,3.24) {};
    \node[bad] at (4.92,2.55) {};
    \node[bad] at (5.25,3.02) {};
    \coordinate (mug) at (1.7,1.3);
    \coordinate (mub) at (4.75,2.92);
    \fill[blue!70!black] (mug) circle (2pt) node[below=2pt] {$\mu_g$};
    \fill[red!70!black] (mub) circle (2pt) node[above=2pt] {$\mu_b$};
    \draw[flow, purple!75!black] (mug) -- node[above, sloped] {$r=\mu_b-\mu_g$} (mub);
    \draw[dashed, black!55] (0.7,2.75) -- (5.9,0.45) node[below left] {$r^\perp$};
    \draw[thinflow, red!70!black] (4.25,2.88) -- (3.37,1.85);
    \draw[thinflow, red!70!black] (4.92,2.55) -- (3.70,1.43);
    \node[align=center] at (3.3,0.32) {删除一个全局分量：$x\mapsto\Pi_r(x)$};
  \end{scope}

  % Panel B: SOM directions
  \begin{scope}[xshift=7.05cm]
    \draw[panel] (0,0) rectangle (7.35,4.65);
    \node[font=\sffamily\bfseries\small] at (3.675,4.34) {(b) SOM 多方向};
    \draw[step=0.72cm, black!35] (0.85,0.75) grid (3.73,3.63);
    \foreach \i in {0,...,3}{
      \foreach \j in {0,...,3}{
        \node[proto] at ({0.85+0.72*\i},{0.75+0.72*\j}) {};
      }
    }
    \fill[blue!70!black] (5.80,1.10) circle (2.2pt) node[below=2pt] {$\nu$};
    \draw[thinflow, orange!80!black] (5.80,1.10) -- node[below, sloped] {$r_1$} (1.57,1.47);
    \draw[thinflow, orange!80!black] (5.80,1.10) -- node[above, sloped] {$r_2$} (3.01,2.91);
    \draw[thinflow, orange!80!black] (5.80,1.10) -- node[above, sloped] {$r_3$} (1.57,3.63);
    \node[align=center] at (5.78,2.55) {$r_i=w_i-\nu$\\16 个候选};
    \node[align=center] at (3.675,0.30) {BO/TPE 选择有序的 $k=2,\ldots,7$ 个投影};
  \end{scope}

  % Panel C: local matrix mapping
  \begin{scope}[xshift=14.85cm]
    \draw[panel] (0,0) rectangle (8.1,4.65);
    \node[font=\sffamily\bfseries\small] at (4.05,4.34) {(c) 上游局部矩阵映射};
    \node[bad] (xb1) at (0.75,2.75) {};
    \node[bad] (xb2) at (0.75,2.15) {};
    \node[bad] (xb3) at (0.75,1.55) {};
    \node[align=center] at (0.75,0.93) {模块输入\\$X_b$};
    \node[draw, rounded corners=2pt, fill=purple!12, minimum width=1.05cm, minimum height=1.05cm] (W) at (2.25,2.15) {$W$};
    \draw[flow] (1.02,2.15) -- (W);
    \node[mapped] (z1) at (4.05,2.80) {};
    \node[mapped] (z2) at (4.38,2.30) {};
    \node[mapped] (z3) at (3.82,1.82) {};
    \draw[flow, purple!75!black] (W) -- node[above] {$X_bW^\top$} (3.62,2.30);
    \node[good] at (5.15,2.92) {};
    \node[good] at (5.52,2.48) {};
    \node[good] at (5.12,1.95) {};
    \node[good] at (5.72,2.02) {};
    \node[align=center, blue!70!black] at (5.45,3.52) {良性输出邻域 $Y_g$};
    \node[bad] at (7.15,1.10) {};
    \node[bad] at (7.50,1.48) {};
    \node[bad] at (7.20,1.86) {};
    \node[align=center, red!70!black] at (7.18,0.48) {原目标输出 $Y_b$};
    \draw[thinflow, green!50!black] (4.38,2.30) -- node[above, sloped] {拉近} (5.12,2.48);
    \draw[thinflow, red!70!black, dashed] (4.05,1.82) -- node[below, sloped] {推离} (6.98,1.30);
    \node[align=center] at (3.95,0.36) {直接优化映射；不显式构造 $r_i$};
  \end{scope}

  \node[font=\sffamily\footnotesize, align=center] at (11.45,-0.52)
    {三种参数化的结构并列；面板顺序不表示理论包含关系，也不表示效果递增。};
  \end{tikzpicture}%
  }
  \caption{从显式方向到直接映射的几何对照。蓝色表示良性参照，红色表示目标行为样本，橙色表示 SOM 原型，绿色方块表示矩阵映射后的代理输出。SOM-MD 仍以方向和投影为显式变量；面板 (c) 示意上游矩阵代理；本地 point-v1 使用缓存输出加低秩增量与软邻域。该图只比较结构，不表示实测轨迹、机制因果性或跨协议性能。}
  \label{fig:method-geometry}
\end{figure*}


单方向方法用两类均值差定义一个全局轴，再删除沿该轴的分量\cite{arditi2024refusal}。SOM-MD 先以 16 个原型近似有害状态的局部结构，再用同一个无害质心生成方向池；BO/TPE 选择 2--7 个方向的排列，最终干预仍由投影组成\cite{piras2026som}。上游矩阵法则把每个目标模块的 $X\mapsto Y$ 映射作为优化对象，其局部 $k$NN 损失约束目标行为新输出与 $Y_g,Y_b$ 的相对距离。这里的“局部邻域映射”不是最优传输：源码没有运输计划、质量守恒或样本间全局耦合。

三种几何锚点保留的信息不同。均值差只保留两组状态的一阶中心；SOM 用 16 个相互关联的原型覆盖有害状态的高密度区域，但仍以单个 $\nu$ 表示无害参照；矩阵法不先聚合 $Y_g,Y_b$，而在每次闭包中重新计算候选输出到参考样本的最近邻。本地 point-v1 将硬近邻改为全参考集上的软邻域，但同样依赖冻结缓存。两种版本保留局部样本结构，却把结果绑定到冻结参考集合和欧氏距离。它不是对“拒答流形”的无条件恢复，也不能证明参考集合之外的状态会按相同方式映射。

干预算子的代数性质也不能只按参数数量排序。单方向投影是固定的低维删除；SOM-MD 复合多个一般不交换的投影；全矩阵能够表示投影、旋转和缩放等线性变化，但源码的逐行范数处理与局部损失会限制实际解。参数空间较大意味着搜索自由度和优化风险同时增加，不意味着矩阵法在因果控制上严格强于方向法。

\subsection{结构比较矩阵}

表~\ref{tab:method-comparison} 用共同坐标区分方向法、上游矩阵原型和本地 point-v1。该表不填写 ASR 或能力分数，因为现有来源没有统一模型、数据划分、判别器、解码配置和计算预算。

\begin{table*}[t]
\centering
\caption{方向法、上游原型与本地 point-v1 的结构比较。SOM-MD 为本文简称；各行没有共享效果评测协议。}
\label{tab:method-comparison}
\footnotesize
\setlength{\tabcolsep}{3.5pt}
\renewcommand{\arraystretch}{1.16}
\begin{tabularx}{\textwidth}{@{}>{\raggedright\arraybackslash}p{21mm}>{\raggedright\arraybackslash}p{24mm}Y Y Y Y@{}}
\toprule
方法 & 决策变量 & 几何锚点 & 干预算子 & 搜索对象 & 证据边界 \\
\midrule
单方向 & 一个方向 $r$ & 有害与无害状态均值差 & 单个正交投影，可折叠进权重 & 层与候选方向筛选 & 同行评议结果限于原论文协议 \\
SOM-MD & 有序方向组 & 16 个有害 SOM 原型减一个无害质心 & 2--7 个投影顺序复合 & BO/TPE 搜索排列 & 原文 7+1 个目标设置；验证/测试关系不确定 \\
上游 ARA & 模块 $W$ 或因子 $A,B$ & 冻结模块 I/O 与硬 $k$NN & 全矩阵写回或低秩适配；可有行归一化 & L-BFGS 与历史分段评分 TPE & 源码形式化与数学审计；无本文归档实测 \\
本地 point-v1 & 秩至多 128 的 $BA$ & 冻结 I/O、归一化软邻域 & 缓存输出加增量；无行归一化 & 六维 TPE，原始 Keywords/KL & 一次 120 次试验验证搜索；联合门槛失败 \\
\bottomrule
\end{tabularx}
\end{table*}


“任意秩”只描述上游全矩阵路径没有显式低秩因子。它不保证 $\Delta W=W-W_0$ 达到任意预指定秩，也不保证满秩。可选低秩路径反而满足 $\operatorname{rank}(BA)\leq r$。同理，矩阵具有更多自由参数并不能单独推出更强的行为控制；冻结 I/O、行范数处理、损失权重和优化轨迹都会限制实际更新。

上游低秩工程路径还暴露了比较表无法承载的实现边界。完整行归一化开启时，优化闭包评价归一化后的 $W_0+BA$，而适配器部署仍使用未投影的 $W_0+BA$；两者一般不同。这个静态差异应被披露，但其对拒答与能力指标的影响尚无实验，不能把它升级为“低秩版本失效”的结论。全矩阵路径在结束时写回归一化矩阵，不存在同一写回差异。本地 point-v1 关闭行归一化并直接在缓存输出上加增量，另用 QR/SVD 检查因子规范化前后的校准输出；因此该历史差异不能用来解释本地失败结果。

\subsection{搜索成本与证据不可比性}

三类方法的外层成本衡量不同对象。单方向流程需要选层和筛选方向；SOM-MD 对有序方向组合执行完整生成与判别器评分，其候选空间随 $k$ 快速增长；矩阵法在每个外层试验中逐模块运行连续 L-BFGS，再用质量与拒答代理评分。即便两套流程都使用 TPE，试验次数也不是等价预算，因为一次试验包含的模块优化、生成数量和判别器调用不同。

证据等级同样不能横向替换。Arditi 等、Wollschl\"ager 等和 Piras 等给出了各自协议内的同行评议结果\cite{arditi2024refusal,wollschlager2025geometry,piras2026som}；上游矩阵原型具有固定源码、数学推论和公开讨论；本地 point-v1 新增一次验收失败的归档搜索，但仍没有同协议受控比较。Piras 等报告的 SOM-MD 结果只能写为“该文在其搜索与评测协议下报告”，还需保留 HarmBench 验证/测试关系未公开的边界。本文据此比较机制，不判断哪一种方法在数值上占优。

公平的后续比较至少需要锁定五组变量。三种方法应使用同一模型版本和目标模块范围，训练、搜索与测试提示应按样本标识公开隔离，生成应共享系统提示和解码参数，行为判别器与无害能力指标应一致，总预算还应同时报告生成次数、优化器闭包调用、GPU 时长和峰值内存。缺少任一组控制时，ASR 差异可能来自搜索预算、数据选择或判别器，而不是几何参数化。
